\documentclass[10pt,a4paper]{amsart}
\usepackage[margin=19mm]{geometry}
\usepackage{amsmath,amssymb,mathtools}
\usepackage[scr=rsfs,cal=euler]{mathalfa}
\usepackage{xfrac}
\usepackage[unicode,hidelinks,bookmarks=false]{hyperref}
\newcommand{\ie}{i.e.\ }
\newcommand{\cf}{cf.\ }
\newcommand{\resp}{respectively\ }
\newcommand{\Subsection}{\subsection}
\providecommand{\nobreakdash}{\nobreak\hbox{-}}
\setlength{\emergencystretch}{2em}
\multlinegap0pt
\usepackage{bm}
\usepackage{bbm}
\usepackage{dsfont}
\usepackage{xspace}
\usepackage{cancel}
\usepackage{mathtools}
\usepackage{amsthm}
\makeatletter
\@ifundefined{lemma}{\newtheorem{lemma}{Lemma}}{}
\makeatother
\title[Notes on the Invariance of Tautness Under Lie Sphere Transfo]{Notes on the Invariance of Tautness Under Lie Sphere Transformations}
\author{Thomas E. Cecil}
\date{Source: arXiv:2506.08834v1 (CC BY 4.0)}
\begin{document}
\maketitle

\noindent\emph{Source and licence.} Notes on the Invariance of Tautness Under Lie Sphere Transformations. Version arXiv:2506.08834v1 (CC BY 4.0), distributed under the Creative Commons Attribution 4.0 International licence, as recorded in the arXiv licence metadata for this exact version and shown on the abstract page https://arxiv.org/abs/2506.08834v1. Full rights notice: licenses/arxiv-2506.08834-CC-BY-4.0.txt.\par\smallskip

\noindent\textit{Editorial scope.} Counted unit: the Lemma whose source label is \texttt{lem:3.6.3} in the article (source0.tex lines 544--555, source label \texttt{lem:3.6.3}), delivered together with the article's own complete original proof of it (source0.tex lines 556--751, one \texttt{proof} environment of 196 lines). No mathematics is written, repaired or re-derived: statement and proof are the article's own text line for line. Required context retained uncounted: eq:3.6.3 (lines 504--507). Every definition, display and label of those lines is retained unchanged. Omitted from this package and not redistributed: 1--503, 508--543, 752--1278. Nothing inside a retained excerpt is omitted or altered: every retained body is delivered exactly as the article's lines are written, so no omission receipt of any kind accompanies this package. Citations: the retained excerpts carry no citation command at all, so no bibliography entry of the article is transcribed and no reference is redistributed. Reference adaptation: none was needed and none was made; every reference inside the delivered unit resolves inside it, so no unresolved reference or citation is produced. Printed slips kept literal: no printed word, symbol, hypothesis or number of the counted statement or of its proof was altered. Layout and class: the article's own class is \texttt{article} with packages including amsmath, amsthm, amssymb, euscript, verbatim; the wrapper is the lane's pinned \texttt{amsart} editorial wrapper at 10pt on A4, which restates the theorem-like environments the retained text needs and adds the article's own macro definitions that the retained text needs (none), with extra wrapper packages (bm, bbm, dsfont, xspace, cancel, mathtools). The delivered numbering is the unit's own. Assets: no retained span contains a figure environment, a graphics-inclusion command, a PDF-page inclusion, a table or a TikZ picture; no image file is copied into this package, the package declares no asset dependency and no image file is redistributed with it. Licence: version arXiv:2506.08834v1 of this article is distributed under the Creative Commons Attribution 4.0 International licence, as recorded in the arXiv licence metadata for this exact version and shown on the abstract page https://arxiv.org/abs/2506.08834v1. A full rights notice is delivered at licenses/arxiv-2506.08834-CC-BY-4.0.txt. Characters: every retained excerpt is pure ASCII, so the delivered engine, XeLaTeX with its default Latin Modern text font, reports no missing character.
\begin{equation}
\label{eq:3.6.3}
\cos r = x \cdot (\cos r\  p + \sin r\  \xi).
\end{equation}

\begin{lemma}
\label{lem:3.6.3} 
Let $\phi:V \rightarrow S^n$ be an immersion of a connected manifold $V$ with $\dim V < n$ into $S^n$, 
and let $(p, \xi) \in T_1S^n$
such that $p \notin \phi (V)$.\\ 
{\rm (a)} A point $x_0 \in V$ is a critical point of the function $r_{(p,\xi)}$ if and only if the sphere
$S_{x_0}$ containing $\phi(x_0)$ in the parabolic pencil of unoriented spheres
determined by $(p,\xi)$ and the submanifold $\phi (V)$ are
tangent at $\phi (x_0)$.\\
{\rm (b)} If $r_{(p,\xi)}$ has a critical point at $x_0 \in V$, then this critical point is degenerate if and only
if the sphere $S_{x_0}$ is a curvature sphere of $\phi (V)$ at $x_0$.
\end{lemma}

\begin{proof}
(a) For any point $x_0 \in V$, there exists a sufficiently small neighborhood $U$ of $x_0$ such that the
restriction of $\phi$ to $U$ is an embedding. We will identify $U$ with its embedded image
$\phi (U) \subset S^n$ and omit the reference to the embedding $\phi$. For brevity, we will denote the
function $r_{(p,\xi)}$ simply as $r$.

Let $X$ be a
smooth vector field tangent to $V$ in such a neighborhood $U$ of $x_0$.  Then on $U$ we compute the derivative
\begin{equation}
\label{eq:3.6.5}
X(\cos r) = - \sin r \ X(r).
\end{equation}
Since $\sin r \neq 0$ for $r \in (0, \pi)$, we see that the functions $\cos r$ and $r$ have the same critical
points on $U$.  Using this and equation \eqref{eq:3.6.3}, we see that
\begin{equation}
\label{eq:3.6.6}
X(r) = 0 \Leftrightarrow X(\cos r) = 0 \Leftrightarrow X(x \cdot (\cos r\  p + \sin r\  \xi)) = 0.
\end{equation}
Then we compute
\begin{equation}
\label{eq:3.6.7}
X(x \cdot (\cos r p + \sin r \xi)) = X \cdot (\cos r p + \sin r \xi) + x \cdot (X(\cos r) p + X(\sin r) \xi).
\end{equation}
If $X(r) = 0$, then $X(\cos r) = X(\sin r) = 0$, and we get from the equation \eqref{eq:3.6.7} above that
\begin{equation}
\label{eq:3.6.8}
X \cdot (\cos r\  p + \sin r\  \xi) = 0.
\end{equation}
Thus, if the function $r = r_{(p,\xi)}$ has a critical point at $x_0$, we have
\begin{equation}
\label{eq:3.6.9}
X \cdot q = 0,
\end{equation}
for all $X \in T_{x_0}V$, where $ q = \cos r\  p + \sin r\  \xi$ is the center of the sphere $S_{x_0}$ in the parabolic
pencil of spheres determined by $(p, \xi)$ that contains the point $x_0$.  The normal space 
to the sphere $S_{x_0}$ in
${\bf R}^{n+1}$ at the point $x_0$
is spanned by the vectors $x_0$ and $q$.  So equation \eqref{eq:3.6.9}, together with the fact that
$X \cdot x_0 = 0$ for $X \in T_{x_0}V$, implies that the tangent space $T_{x_0}V$ is contained in the tangent
space $T_{x_0}S_{x_0}$, i.e., the sphere $S_{x_0}$ and the submanifold $\phi (V)$ are tangent at the point $\phi (x_0)$.

Conversely, if the tangent space $T_{x_0}V$ is contained in the tangent
space $T_{x_0}S_{x_0}$, then $X(r) = 0$ for all $X \in T_{x_0}V$, because $S_{x_0} - \{p\}$ is a level set of the function
$r$ in $S^n$. Thus, $r$ has a critical point at $x_0$.\\

\noindent 
(b) We now want to compute the 
Hessian of the function $r$ at a critical point $x_0 \in V$.  Let $X$ and $Y$ be
smooth vector fields tangent to $V$ on the neighborhood $U$ of $x_0$ in $V$ that was used in part (a).  
Then the value $H(X,Y)$ of 
the Hessian of $r$
at $x_0$ equals $YX(r)$ at $x_0$.  We first note that
\begin{equation}
\label{eq:3.6.10}
YX(\cos r) = Y(-\sin r \ X(r)) = Y(-\sin r) X(r) - \sin r \ YX(r).
\end{equation}
At the critical point $x_0$, we have $X(r) = 0$, and thus
\begin{equation}
\label{eq:3.6.11}
YX(\cos r) = - \sin r \ YX(r).
\end{equation}
On the other hand, by differentiating equation \eqref{eq:3.6.3} we get
\begin{equation}
\label{eq:3.6.12}
X(\cos r) = X \cdot (\cos r\  p + \sin r\  \xi) + x \cdot (X(\cos r) p + X(\sin r) \xi).
\end{equation}
Then by differentiating equation \eqref{eq:3.6.12}, we have
\begin{eqnarray}
\label{eq:3.6.13}
YX(\cos r) = D_Y X \cdot (\cos r\  p + \sin r\  \xi) + X \cdot (Y(\cos r) p + Y(\sin r) \xi) \\ 
 + \ Y \cdot (X(\cos r) p + X(\sin r) \xi) + x \cdot (YX(\cos r) p + YX(\sin r) \xi), \nonumber
\end{eqnarray}
where $D$ is the Euclidean covariant derivative on ${\bf R}^{n+1}$.  At the critical point $x_0$, we have
\begin{equation}
\label{eq:3.6.14}
X(\cos r) = X(\sin r) = Y(\cos r) = Y(\sin r) = 0,
\end{equation}
and so equation \eqref{eq:3.6.13} becomes
\begin{equation}
\label{eq:3.6.15}
YX(\cos r) = D_Y X \cdot (\cos r\  p + \sin r\  \xi) + x_0 \cdot (YX(\cos r) p + YX(\sin r) \xi).
\end{equation}
We now examine the two terms on the right side of equation \eqref{eq:3.6.15} separately.  Note that at the critical point
$x_0$, the vector $q = \cos r \ p + \sin r \ \xi$ lies in the normal
space to $\phi (V)$ at $x_0$, and so
\begin{equation}
\label{eq:3.6.16}
q = \cos r\  x_0 + \sin r\  N, 
\end{equation}
where $N$ is a unit normal vector to $\phi (V)$ at $x_0$.  Thus we can write the covariant derivative
$D_Y X$ as
\begin{equation}
\label{eq:3.6.17}
D_Y X = \nabla_Y X + \alpha (X,Y) - (X \cdot Y) x_0, 
\end{equation}
where $\nabla$ is the Levi--Civita connection on $\phi (V)$ induced from $D$, $\alpha$ is the second fundamental
form of $\phi (V)$ in $S^n$, and the term $- (X \cdot Y) x_0$ is the second fundamental
form of the sphere $S^n$ as a hypersurface in ${\bf R}^{n+1}$. 
Then using equation \eqref{eq:3.6.16} we obtain
\begin{eqnarray}
\label{eq:3.6.18}
D_Y X \cdot q &=& (\nabla_Y X + \alpha (X,Y) - (X \cdot Y) x_0) \cdot q \\
 &=& \alpha (X,Y) \cdot (\cos r\  x_0 + \sin r\  N) - (X \cdot Y) x_0 \cdot (\cos r\  x_0 + \sin r\  N) \nonumber \\
 &=& \sin r \ A_N X \cdot Y - \cos r \ X \cdot Y, \nonumber
\end{eqnarray}
since $\nabla_Y X$ is orthogonal to $q$, $\alpha (X,Y) \cdot x_0 = 0$, $N \cdot x_0 = 0$,  and 
\begin{displaymath}
\alpha (X,Y)\cdot N = A_N X \cdot Y,
\end{displaymath}
where $A_N$ is the shape operator
determined by the unit normal $N$ to $\phi (V)$ at $x_0$.

Next we consider the second term on the right side of equation \eqref{eq:3.6.15},
\begin{equation}
\label{eq:3.6.19}
x_0 \cdot (YX(\cos r) p + YX(\sin r) \xi).
\end{equation}
Note that at the critical point $x_0$ we have $X(r) = Y(r) = 0$, and so
\begin{eqnarray}
\label{eq:3.6.20}
YX(\cos r) &=& Y( - \sin r \ X(r)) = Y(- \sin r) X(r) - \sin r \ YX(r)\\ \nonumber 
&=& - \sin r \ YX(r),
\end{eqnarray}
and similarly
\begin{equation}
\label{eq:3.6.21}
YX(\sin r) = \cos r \ YX(r).
\end{equation}
Thus we have
\begin{eqnarray}
\label{eq:3.6.22}
YX (\cos r) p + YX(\sin r) \xi &=& - \sin r \ YX(r) p + \cos r \ YX(r) \xi \\
 &=& YX(r) ( -\sin r\  p + \cos r\  \xi). \nonumber 
\end{eqnarray}
So the term in equation \eqref{eq:3.6.19} above is
\begin{equation}
\label{eq:3.6.23}
x_0 \cdot (YX(\cos r) p + YX(\sin r) \xi) = YX(r) (x_0 \cdot ( -\sin r\  p + \cos r\  \xi)).
\end{equation}
From equations \eqref{eq:3.6.15}, \eqref{eq:3.6.18} and \eqref{eq:3.6.23}, we have
\begin{equation}
\label{eq:3.6.24}
YX(\cos r) = \sin r A_N X \cdot Y - \cos r X \cdot Y + YX(r) (x_0 \cdot ( -\sin r\  p + \cos r\  \xi)).
\end{equation}
Using equations \eqref{eq:3.6.11} and \eqref{eq:3.6.24}, we get
\begin{equation}
\label{eq:3.6.25}
(- \sin r - (x_0 \cdot ( -\sin r\  p + \cos r\  \xi))) YX(r) = \sin r A_N X \cdot Y - \cos r X \cdot Y.
\end{equation}
Denote the coefficient of $YX(r)$ in equation \eqref{eq:3.6.25} by 
\begin{equation}
\label{eq:3.6.26}
C = - \sin r - (x_0 \cdot ( -\sin r\  p + \cos r\  \xi)).
\end{equation}
Note that the vector 
\begin{equation}
\label{eq:3.6.27}
v =  -\sin r\  p + \cos r\  \xi
\end{equation}
is tangent at the point $q$ to the geodesic in $S^n$ from $p$ with initial tangent vector $\xi$.



On the sphere
$S_{x_0}$ centered at $q$, the function $x \cdot v$ takes its minimum value at $p$ where it is equal to $- \sin r$.  Thus
\begin{equation}
\label{eq:3.6.28}
x_0 \cdot ( -\sin r\  p + \cos r\  \xi) > - \sin r,
\end{equation}
since $x_0 \in S_{x_0}$ and $x_0 \neq p$.  So the term
\begin{equation}
\label{eq:3.6.29}
\sin r + x_0 \cdot ( -\sin r\  p + \cos r\  \xi) > 0,
\end{equation}
and the coefficient $C$ in equation \eqref{eq:3.6.26} is negative.  Thus we have from equation \eqref{eq:3.6.25},
\begin{eqnarray}
\label{eq:3.6.30}
YX(r) &=& (1/C) (\sin r \ (A_N X \cdot Y) - \cos r \ (X \cdot Y))\\
 &=& (1/C) (\sin r \ A_N - \cos r \ I) X \cdot Y.\nonumber
\end{eqnarray}
Thus the Hessian $H(X,Y)$ of $r$ at $x_0$ is degenerate if and only if there is a nonzero vector $X \in T_{x_0} V$
such that 
\begin{equation}
\label{eq:3.6.31}
(\sin r \ A_N - \cos r \ I) X = 0,
\end{equation}
that is,
\begin{equation}
\label{eq:3.6.32}
A_N X = \cot r \ X.
\end{equation}
This equation holds if and only if
$\cot r$ is an eigenvalue of $A_N$ with corresponding principal vector $X$, i.e., the point
$q = \cos r\  x_0 + \sin r\  N$ is a focal point of $\phi (V)$ at $x_0$, and thus the corresponding sphere $S_{x_0}$ is
a curvature sphere of $\phi (V)$ at $x_0$. 
\end{proof}

\end{document}
